\scanpage{32}
\mspage{30}
% archive sheet number: 31
% hand: Grothendieck
% sheet_kind: text
% element: text | upper left | I. 6. Prop. 6.1. Soit
I

6.

\uline{Prop.} 6.1. Soit
% element: display | top, line 4 | 1 I pi pi0 1
\[
1\longrightarrow I\unsure{-|\longrightarrow}\pi-\pi_0\longrightarrow1
\]
% element: text | upper, lines 5-10 | une extension de groupes profinis
une extension de groupes profinis, et supposons que pour \unsure{tt.|tt|$\forall$} s-groupe ouvert \deleted{\unsure{le}} $\pi'$ de $\pi$, posant $I'=\deleted{\unsure{\pi}}I\cap\pi'$, \unsure{le} \deleted{\unsure{l'action de} $\pi'$ \unsure{sur}} \deleted{groupe $[I',\pi']$ est dense dans $I'$, i.e. le groupe $I'^{\mathrm{ab}}(\ell)_{\deleted{\unsure{\pi_0}}}$}, plus grand quotient profini \inserted{$(I'^{\mathrm{ab}}(\ell))_{\pi_0}$} de $I'^{\mathrm{ab}}(\ell)$ sur lequel $\pi'_0=\pi'/I'$ opère trivialement, est un groupe fini. Soit
% element: display | upper-middle, line 11 | u pi AutQell(Eell)
\[
u:\pi\longrightarrow\operatorname{Aut}_{\mathbb Q_\ell}(E_\ell)
\]
% element: text | middle, lines 12-14 | une représentation ell-adique
une représentation $\ell$-adique de $\pi$, \deleted{\unsure{soit et}} considérons \deleted{la représentation induite par sur $I$} ss-groupe algébrique \inserted{$H$ ($G$)} de $\operatorname{Aut}_{\mathbb Q_\ell}(E_\ell)$ \uline{engendré} par $u(I)_{\unsure{k|\ell}}$ (resp. $\unsure{I(G)|u(\pi)}$). Alors

% element: text | middle, lines 15-16 | a) H est invariant dans G
 a) $H$ est invariant dans $G$, donc $H^\circ$ invariant dans $G$ \deleted{et \unsure{a fortiori}, dans $G^\circ$}.

% element: text | middle-lower, lines 17-20 | b) Soit R le radical de H0
 b) \deleted{\unsure{L'opération} de $G^\circ$ sur} Soit $R$ le radical de $H^\circ$, \inserted{(qui est inv. dans $G$, et soit \deleted{$\unsure{A\subset R}$} $A=R/[R,R]$)} Alors l'action de $G^\circ$ sur $A$ \deleted{\unsure{identifie}} \inserted{satisfait} \deleted{\unsure{est intérieure et|est trivialement et} \unsure{l'identité sur}} la condition $A_{(G^\circ)}=0$, i.e. $A$ est \deleted{\unsure{le groupe engendré}} \uline{engendré} par les $(1-g)(A)$, où $g\in G^\circ(\mathbb Q_\ell)$.

% element: text | lower, line 21 | c) R est un groupe unipotent
 c) $R$ est un groupe \uline{unipotent}.

% element: text | lower, lines 22-24 | Dém a) Résulte trivialt
\uline{Dém} a) Résulte \unsure{trivialt|trivialement} du fait que $I$ inv. dans $\pi$. Il s'ensuit que $R$ est aussi inv. dans $G$, donc $G$ opère sur $A=R/[R,R]$.

% element: text | bottom, lines 25-26 | b) Exprimer l'hypothèse
 b) Exprimer l'hypothèse $I'_{\mathrm{ab}}(\ell)$ fini, en notant que $H^\circ/[H^\circ,H^\circ]\longleftarrow R/[R,R]=A$ \unsure{est} une \uline{isogénie} \deleted{\unsure{propre}}.

% element: text | bottom, lines 27-30 | c) A priori R=T.U
 c) A priori, $R=T.U$, $U$ unipotent, $T$ tore ($U$ ss-groupe car. de $R$). \deleted{\unsure{Et}} Le groupe $U$ est invariant \deleted{dans $R$, et $G$ opère} sur $R/U=T$, de façon que $G^\circ$ opère trivialement (car les aut. d'un tore forment
